\documentclass[12pt, letterpaper]{article}

\usepackage[utf8]{inputenc}
\usepackage[T1]{fontenc}
\usepackage[spanish]{babel}
\usepackage{geometry}
\usepackage{amsmath, amssymb, amsthm}
\usepackage{booktabs}
\usepackage{array}
\usepackage{microtype}
\usepackage{setspace}
\usepackage{parskip}
\usepackage{titlesec}
\usepackage{fancyhdr}
\usepackage{enumitem}
\usepackage{bm}
\usepackage{mathtools}
\usepackage{tcolorbox}
\usepackage{hyperref}
\usepackage{pgfplots}
\usepackage{pgfplotstable}
\usepackage{xcolor}
\usepackage{colortbl}
\usepackage{caption}

\pgfplotsset{compat=1.18}

\tcbuselibrary{skins, breakable}

\geometry{top=2.8cm, bottom=3.0cm, left=3.2cm, right=3.2cm}
\setstretch{1.20}

\pagestyle{fancy}
\fancyhf{}
\fancyhead[C]{\small\textsc{Econometr\'ia --- Gujarati \& Porter --- Ejercicio 1.2}}
\fancyfoot[C]{\small\thepage}
\renewcommand{\headrulewidth}{0.4pt}
\renewcommand{\footrulewidth}{0pt}

\titleformat{\section}
  {\large\bfseries\scshape}{\thesection.}{0.6em}{}
  [\vspace{-4pt}\rule{\linewidth}{0.4pt}]
\titleformat{\subsection}{\normalsize\bfseries}{\thesubsection.}{0.5em}{}
\titleformat{\subsubsection}{\normalsize\itshape}{}{0em}{}

\newtheorem{prop}{Proposici\'on}
\theoremstyle{remark}
\newtheorem{obs}{Observaci\'on}

\newtcolorbox{clave}{
  colback=white, colframe=black,
  boxrule=0.5pt, arc=3pt,
  left=8pt, right=8pt, top=6pt, bottom=6pt,
  breakable
}

\newtcolorbox{datebox}{
  colback=gray!4, colframe=black!30,
  boxrule=0.4pt, arc=2pt,
  left=8pt, right=8pt, top=5pt, bottom=5pt,
  breakable
}

% Colores por pais
\definecolor{col_usa}{HTML}{2c3e50}
\definecolor{col_canada}{HTML}{1b6ca8}
\definecolor{col_japan}{HTML}{d35400}
\definecolor{col_france}{HTML}{c0392b}
\definecolor{col_germany}{HTML}{27ae60}
\definecolor{col_italy}{HTML}{8e44ad}
\definecolor{col_uk}{HTML}{16a085}

\newcommand{\Var}{\operatorname{Var}}
\newcommand{\Corr}{\operatorname{Corr}}

\begin{document}

% ---------------------------------------------------------------
%  Portada
% ---------------------------------------------------------------
\begin{center}
  {\LARGE\bfseries Comparaci\'on de Tasas de Inflaci\'on}\\[6pt]
  {\LARGE\bfseries frente a Estados Unidos (1981--2005)}\\[12pt]
  {\large\itshape Gr\'aficos comparativos, an\'alisis de co-movimiento}\\
  {\large\itshape y discusi\'on sobre causalidad vs.\ correlaci\'on}\\[14pt]
  \rule{0.55\linewidth}{0.6pt}\\[8pt]
  {\normalsize Ejercicio~1.2 --- \textit{Econometr\'ia}, Gujarati \& Porter, 5.a ed.}\\[4pt]
  {\small\textsc{Fuente: Tabla~1.3 --- IPC base 1982--1984 $=$ 100, 1980--2005}}
\end{center}

\vspace{10pt}

\noindent\textbf{El ejercicio.}\quad
La Tabla~1.3 reporta el IPC de siete pa\'ises para el per\'iodo 1980--2005.
Se solicita: \emph{(a)} graficar la tasa de inflaci\'on de los seis pa\'ises
en comparaci\'on con la de EE.\,UU.; \emph{(b)} comentar el comportamiento
comparativo; y \emph{(c)} discutir si la co-direcci\'on de las tasas implica
que EE.\,UU.\ \emph{causa} la inflaci\'on en los dem\'as pa\'ises.

% ---------------------------------------------------------------
\section{C\'alculo de tasas de inflaci\'on}
% ---------------------------------------------------------------

La tasa de inflaci\'on se calcula como:
\begin{equation}
  \pi_t = \frac{\text{IPC}_t - \text{IPC}_{t-1}}{\text{IPC}_{t-1}} \times 100
  \label{eq:pi}
\end{equation}
Por ejemplo, para Canad\'a en 1981:
\begin{equation}
  \pi_{1981}^{\text{CA}} = \frac{85.6 - 76.1}{76.1}\times 100
  = \frac{9.5}{76.1}\times 100 \approx 12.48\%
  \label{eq:ejemplo}
\end{equation}

\subsection{Tasas de inflaci\'on calculadas (\%)}

\begin{datebox}
\small\centering
\captionsetup{type=table}
\captionof{table}{Tasas de inflaci\'on anual (\%) calculadas con \eqref{eq:pi},
  1981--2005}
\vspace{4pt}
\begin{tabular}{crrrrrrrr}
\toprule
\textbf{A\~no} & \textbf{EE.\,UU.} & \textbf{Canad\'a} & \textbf{Jap\'on} & \textbf{Francia} & \textbf{Alemania} & \textbf{Italia} & \textbf{G.B.} \\
\midrule
1981 & 10.32 & 12.48 &  4.73 & 13.30 &  6.34 & 18.15 & 11.97 \\
1982 &  6.16 & 10.86 &  2.94 & 12.10 &  5.20 & 16.29 &  8.54 \\
1983 &  3.21 &  5.80 &  1.73 &  9.37 &  3.41 & 14.91 &  4.61 \\
1984 &  4.32 &  4.28 &  2.30 &  7.68 &  2.39 & 10.52 &  5.01 \\
1985 &  3.56 &  4.11 &  2.06 &  5.83 &  2.04 &  9.24 &  6.01 \\
1986 &  1.86 &  4.13 &  0.67 &  2.53 & -0.19 &  5.91 &  3.42 \\
1987 &  3.65 &  4.32 &  0.00 &  3.33 &  0.29 &  4.81 &  4.18 \\
1988 &  4.14 &  4.05 &  0.67 &  2.64 &  1.34 &  5.03 &  4.93 \\
1989 &  4.82 &  4.95 &  2.27 &  3.54 &  2.73 &  6.20 &  7.80 \\
1990 &  5.40 &  4.80 &  3.15 &  3.26 &  2.75 &  6.44 &  9.47 \\
1991 &  4.21 &  5.61 &  3.23 &  3.23 &  3.65 &  6.30 &  5.87 \\
1992 &  3.01 &  1.54 &  1.74 &  2.33 &  5.07 &  5.28 &  3.57 \\
1993 &  2.99 &  1.79 &  1.28 &  2.14 &  4.42 &  4.57 &  1.60 \\
1994 &  2.56 &  0.20 &  0.67 &  1.67 &  2.74 &  4.05 &  2.42 \\
1995 &  2.83 &  2.16 & -0.08 &  1.78 &  1.68 &  5.27 &  3.48 \\
1996 &  2.95 &  1.58 &  0.08 &  2.02 &  1.50 &  4.08 &  2.40 \\
1997 &  2.29 &  1.62 &  1.84 &  1.19 &  1.85 &  2.05 &  3.18 \\
1998 &  1.56 &  0.96 &  0.58 &  0.65 &  0.95 &  2.89 &  3.40 \\
1999 &  2.21 &  1.71 & -0.33 &  0.52 &  0.65 &  1.66 &  1.51 \\
2000 &  3.37 &  2.74 & -0.66 &  1.68 &  1.43 &  2.52 &  2.98 \\
2001 &  2.84 &  2.55 & -0.74 &  1.65 &  1.97 &  2.76 &  1.75 \\
2002 &  1.58 &  2.25 & -0.92 &  1.93 &  1.31 &  2.52 &  1.67 \\
2003 &  2.28 &  2.78 & -0.25 &  2.08 &  1.09 &  2.66 &  2.90 \\
2004 &  2.66 &  1.86 &  0.00 &  2.16 &  1.68 &  2.19 &  2.99 \\
2005 &  3.39 &  2.15 & -0.34 &  1.70 &  1.92 &  1.95 &  2.83 \\
\midrule
\textbf{Media}  & 3.43 & 3.78 & 0.83 & 3.65 & 2.33 & 6.55 & 4.52 \\
\textbf{Desv.Est.} & 1.80 & 2.79 & 1.42 & 3.37 & 1.80 & 4.39 & 2.68 \\
\bottomrule
\end{tabular}
\end{datebox}

% ---------------------------------------------------------------
\section{Gr\'aficos comparativos frente a EE.\,UU.}
% ---------------------------------------------------------------

Cada gr\'afico superpone la tasa de inflaci\'on de un pa\'is (en color)
sobre la de EE.\,UU.\ (trazo negro discontinuo), permitiendo identificar
visualmente el co-movimiento, las brechas de nivel y las divergencias.

% Macro comun para el eje de EE.UU.
\newcommand{\plotUSA}{
  \addplot[color=col_usa, thick, dashed, mark=none] coordinates {
    (1981,10.32)(1982,6.16)(1983,3.21)(1984,4.32)(1985,3.56)
    (1986,1.86)(1987,3.65)(1988,4.14)(1989,4.82)(1990,5.40)
    (1991,4.21)(1992,3.01)(1993,2.99)(1994,2.56)(1995,2.83)
    (1996,2.95)(1997,2.29)(1998,1.56)(1999,2.21)(2000,3.37)
    (2001,2.84)(2002,1.58)(2003,2.28)(2004,2.66)(2005,3.39)
  };
  \addlegendentry{EE.\,UU.\ (ref.)}
}

% ---------- Canada ----------
\begin{figure}[htbp]
\centering
\begin{tikzpicture}
\begin{axis}[
  width=13cm, height=5.5cm,
  xlabel={A\~no}, ylabel={Inflaci\'on (\%)},
  xmin=1980, xmax=2006, ymin=-2, ymax=22,
  xtick={1981,1985,1990,1995,2000,2005},
  xticklabel style={font=\small},
  yticklabel style={font=\small},
  grid=both,
  grid style={line width=0.2pt, draw=gray!25},
  major grid style={line width=0.35pt, draw=gray!45},
  legend pos=north east,
  legend style={font=\small, draw=black!30, fill=white, inner sep=3pt},
  tick align=outside,
  axis line style={line width=0.5pt},
  title={Canad\'a vs.\ EE.\,UU.}
]
\plotUSA
\addplot[color=col_canada, thick, mark=*, mark size=1.8pt] coordinates {
  (1981,12.48)(1982,10.86)(1983,5.80)(1984,4.28)(1985,4.11)
  (1986,4.13)(1987,4.32)(1988,4.05)(1989,4.95)(1990,4.80)
  (1991,5.61)(1992,1.54)(1993,1.79)(1994,0.20)(1995,2.16)
  (1996,1.58)(1997,1.62)(1998,0.96)(1999,1.71)(2000,2.74)
  (2001,2.55)(2002,2.25)(2003,2.78)(2004,1.86)(2005,2.15)
};
\addlegendentry{Canad\'a}
\end{axis}
\end{tikzpicture}
\caption{Tasa de inflaci\'on: Canad\'a vs.\ EE.\,UU.\ (1981--2005).}
\end{figure}

% ---------- Japon ----------
\begin{figure}[htbp]
\centering
\begin{tikzpicture}
\begin{axis}[
  width=13cm, height=5.5cm,
  xlabel={A\~no}, ylabel={Inflaci\'on (\%)},
  xmin=1980, xmax=2006, ymin=-2, ymax=22,
  xtick={1981,1985,1990,1995,2000,2005},
  xticklabel style={font=\small},
  yticklabel style={font=\small},
  grid=both,
  grid style={line width=0.2pt, draw=gray!25},
  major grid style={line width=0.35pt, draw=gray!45},
  legend pos=north east,
  legend style={font=\small, draw=black!30, fill=white, inner sep=3pt},
  tick align=outside,
  axis line style={line width=0.5pt},
  title={Jap\'on vs.\ EE.\,UU.}
]
\plotUSA
\addplot[color=col_japan, thick, mark=triangle*, mark size=1.8pt] coordinates {
  (1981,4.73)(1982,2.94)(1983,1.73)(1984,2.30)(1985,2.06)
  (1986,0.67)(1987,0.00)(1988,0.67)(1989,2.27)(1990,3.15)
  (1991,3.23)(1992,1.74)(1993,1.28)(1994,0.67)(1995,-0.08)
  (1996,0.08)(1997,1.84)(1998,0.58)(1999,-0.33)(2000,-0.66)
  (2001,-0.74)(2002,-0.92)(2003,-0.25)(2004,0.00)(2005,-0.34)
};
\addlegendentry{Jap\'on}
\end{axis}
\end{tikzpicture}
\caption{Tasa de inflaci\'on: Jap\'on vs.\ EE.\,UU.\ (1981--2005).}
\end{figure}

% ---------- Francia ----------
\begin{figure}[htbp]
\centering
\begin{tikzpicture}
\begin{axis}[
  width=13cm, height=5.5cm,
  xlabel={A\~no}, ylabel={Inflaci\'on (\%)},
  xmin=1980, xmax=2006, ymin=-2, ymax=22,
  xtick={1981,1985,1990,1995,2000,2005},
  xticklabel style={font=\small},
  yticklabel style={font=\small},
  grid=both,
  grid style={line width=0.2pt, draw=gray!25},
  major grid style={line width=0.35pt, draw=gray!45},
  legend pos=north east,
  legend style={font=\small, draw=black!30, fill=white, inner sep=3pt},
  tick align=outside,
  axis line style={line width=0.5pt},
  title={Francia vs.\ EE.\,UU.}
]
\plotUSA
\addplot[color=col_france, thick, mark=square*, mark size=1.8pt] coordinates {
  (1981,13.30)(1982,12.10)(1983,9.37)(1984,7.68)(1985,5.83)
  (1986,2.53)(1987,3.33)(1988,2.64)(1989,3.54)(1990,3.26)
  (1991,3.23)(1992,2.33)(1993,2.14)(1994,1.67)(1995,1.78)
  (1996,2.02)(1997,1.19)(1998,0.65)(1999,0.52)(2000,1.68)
  (2001,1.65)(2002,1.93)(2003,2.08)(2004,2.16)(2005,1.70)
};
\addlegendentry{Francia}
\end{axis}
\end{tikzpicture}
\caption{Tasa de inflaci\'on: Francia vs.\ EE.\,UU.\ (1981--2005).}
\end{figure}

% ---------- Alemania ----------
\begin{figure}[htbp]
\centering
\begin{tikzpicture}
\begin{axis}[
  width=13cm, height=5.5cm,
  xlabel={A\~no}, ylabel={Inflaci\'on (\%)},
  xmin=1980, xmax=2006, ymin=-2, ymax=22,
  xtick={1981,1985,1990,1995,2000,2005},
  xticklabel style={font=\small},
  yticklabel style={font=\small},
  grid=both,
  grid style={line width=0.2pt, draw=gray!25},
  major grid style={line width=0.35pt, draw=gray!45},
  legend pos=north east,
  legend style={font=\small, draw=black!30, fill=white, inner sep=3pt},
  tick align=outside,
  axis line style={line width=0.5pt},
  title={Alemania vs.\ EE.\,UU.}
]
\plotUSA
\addplot[color=col_germany, thick, mark=diamond*, mark size=1.8pt] coordinates {
  (1981,6.34)(1982,5.20)(1983,3.41)(1984,2.39)(1985,2.04)
  (1986,-0.19)(1987,0.29)(1988,1.34)(1989,2.73)(1990,2.75)
  (1991,3.65)(1992,5.07)(1993,4.42)(1994,2.74)(1995,1.68)
  (1996,1.50)(1997,1.85)(1998,0.95)(1999,0.65)(2000,1.43)
  (2001,1.97)(2002,1.31)(2003,1.09)(2004,1.68)(2005,1.92)
};
\addlegendentry{Alemania}
\end{axis}
\end{tikzpicture}
\caption{Tasa de inflaci\'on: Alemania vs.\ EE.\,UU.\ (1981--2005).}
\end{figure}

% ---------- Italia ----------
\begin{figure}[htbp]
\centering
\begin{tikzpicture}
\begin{axis}[
  width=13cm, height=5.5cm,
  xlabel={A\~no}, ylabel={Inflaci\'on (\%)},
  xmin=1980, xmax=2006, ymin=-2, ymax=22,
  xtick={1981,1985,1990,1995,2000,2005},
  xticklabel style={font=\small},
  yticklabel style={font=\small},
  grid=both,
  grid style={line width=0.2pt, draw=gray!25},
  major grid style={line width=0.35pt, draw=gray!45},
  legend pos=north east,
  legend style={font=\small, draw=black!30, fill=white, inner sep=3pt},
  tick align=outside,
  axis line style={line width=0.5pt},
  title={Italia vs.\ EE.\,UU.}
]
\plotUSA
\addplot[color=col_italy, thick, mark=pentagon*, mark size=1.8pt] coordinates {
  (1981,18.15)(1982,16.29)(1983,14.91)(1984,10.52)(1985,9.24)
  (1986,5.91)(1987,4.81)(1988,5.03)(1989,6.20)(1990,6.44)
  (1991,6.30)(1992,5.28)(1993,4.57)(1994,4.05)(1995,5.27)
  (1996,4.08)(1997,2.05)(1998,2.89)(1999,1.66)(2000,2.52)
  (2001,2.76)(2002,2.52)(2003,2.66)(2004,2.19)(2005,1.95)
};
\addlegendentry{Italia}
\end{axis}
\end{tikzpicture}
\caption{Tasa de inflaci\'on: Italia vs.\ EE.\,UU.\ (1981--2005).}
\end{figure}

% ---------- Gran Bretana ----------
\begin{figure}[htbp]
\centering
\begin{tikzpicture}
\begin{axis}[
  width=13cm, height=5.5cm,
  xlabel={A\~no}, ylabel={Inflaci\'on (\%)},
  xmin=1980, xmax=2006, ymin=-2, ymax=22,
  xtick={1981,1985,1990,1995,2000,2005},
  xticklabel style={font=\small},
  yticklabel style={font=\small},
  grid=both,
  grid style={line width=0.2pt, draw=gray!25},
  major grid style={line width=0.35pt, draw=gray!45},
  legend pos=north east,
  legend style={font=\small, draw=black!30, fill=white, inner sep=3pt},
  tick align=outside,
  axis line style={line width=0.5pt},
  title={Gran Breta\~na vs.\ EE.\,UU.}
]
\plotUSA
\addplot[color=col_uk, thick, mark=o, mark size=2pt] coordinates {
  (1981,11.97)(1982,8.54)(1983,4.61)(1984,5.01)(1985,6.01)
  (1986,3.42)(1987,4.18)(1988,4.93)(1989,7.80)(1990,9.47)
  (1991,5.87)(1992,3.57)(1993,1.60)(1994,2.42)(1995,3.48)
  (1996,2.40)(1997,3.18)(1998,3.40)(1999,1.51)(2000,2.98)
  (2001,1.75)(2002,1.67)(2003,2.90)(2004,2.99)(2005,2.83)
};
\addlegendentry{Gran Breta\~na}
\end{axis}
\end{tikzpicture}
\caption{Tasa de inflaci\'on: Gran Breta\~na vs.\ EE.\,UU.\ (1981--2005).}
\end{figure}

% ---------- Comparativo general ----------
\begin{figure}[htbp]
\centering
\begin{tikzpicture}
\begin{axis}[
  width=14cm, height=7cm,
  xlabel={A\~no}, ylabel={Inflaci\'on (\%)},
  xmin=1980, xmax=2006, ymin=-2, ymax=22,
  xtick={1981,1985,1990,1995,2000,2005},
  xticklabel style={font=\small, rotate=45, anchor=east},
  yticklabel style={font=\small},
  grid=both,
  grid style={line width=0.2pt, draw=gray!25},
  major grid style={line width=0.35pt, draw=gray!45},
  legend columns=4,
  legend style={font=\footnotesize, draw=black!30, fill=white,
    inner sep=3pt, at={(0.5,-0.24)}, anchor=north},
  tick align=outside,
  axis line style={line width=0.5pt},
  title={Todos los pa\'ises vs.\ EE.\,UU.\ (1981--2005)}
]
\plotUSA
\addplot[color=col_canada, thick, mark=*, mark size=1.4pt] coordinates {
  (1981,12.48)(1982,10.86)(1983,5.80)(1984,4.28)(1985,4.11)
  (1986,4.13)(1987,4.32)(1988,4.05)(1989,4.95)(1990,4.80)
  (1991,5.61)(1992,1.54)(1993,1.79)(1994,0.20)(1995,2.16)
  (1996,1.58)(1997,1.62)(1998,0.96)(1999,1.71)(2000,2.74)
  (2001,2.55)(2002,2.25)(2003,2.78)(2004,1.86)(2005,2.15)
}; \addlegendentry{Canad\'a}
\addplot[color=col_japan, thick, mark=triangle*, mark size=1.4pt] coordinates {
  (1981,4.73)(1982,2.94)(1983,1.73)(1984,2.30)(1985,2.06)
  (1986,0.67)(1987,0.00)(1988,0.67)(1989,2.27)(1990,3.15)
  (1991,3.23)(1992,1.74)(1993,1.28)(1994,0.67)(1995,-0.08)
  (1996,0.08)(1997,1.84)(1998,0.58)(1999,-0.33)(2000,-0.66)
  (2001,-0.74)(2002,-0.92)(2003,-0.25)(2004,0.00)(2005,-0.34)
}; \addlegendentry{Jap\'on}
\addplot[color=col_france, thick, mark=square*, mark size=1.4pt] coordinates {
  (1981,13.30)(1982,12.10)(1983,9.37)(1984,7.68)(1985,5.83)
  (1986,2.53)(1987,3.33)(1988,2.64)(1989,3.54)(1990,3.26)
  (1991,3.23)(1992,2.33)(1993,2.14)(1994,1.67)(1995,1.78)
  (1996,2.02)(1997,1.19)(1998,0.65)(1999,0.52)(2000,1.68)
  (2001,1.65)(2002,1.93)(2003,2.08)(2004,2.16)(2005,1.70)
}; \addlegendentry{Francia}
\addplot[color=col_germany, thick, mark=diamond*, mark size=1.4pt] coordinates {
  (1981,6.34)(1982,5.20)(1983,3.41)(1984,2.39)(1985,2.04)
  (1986,-0.19)(1987,0.29)(1988,1.34)(1989,2.73)(1990,2.75)
  (1991,3.65)(1992,5.07)(1993,4.42)(1994,2.74)(1995,1.68)
  (1996,1.50)(1997,1.85)(1998,0.95)(1999,0.65)(2000,1.43)
  (2001,1.97)(2002,1.31)(2003,1.09)(2004,1.68)(2005,1.92)
}; \addlegendentry{Alemania}
\addplot[color=col_italy, thick, mark=pentagon*, mark size=1.4pt] coordinates {
  (1981,18.15)(1982,16.29)(1983,14.91)(1984,10.52)(1985,9.24)
  (1986,5.91)(1987,4.81)(1988,5.03)(1989,6.20)(1990,6.44)
  (1991,6.30)(1992,5.28)(1993,4.57)(1994,4.05)(1995,5.27)
  (1996,4.08)(1997,2.05)(1998,2.89)(1999,1.66)(2000,2.52)
  (2001,2.76)(2002,2.52)(2003,2.66)(2004,2.19)(2005,1.95)
}; \addlegendentry{Italia}
\addplot[color=col_uk, thick, mark=o, mark size=1.4pt] coordinates {
  (1981,11.97)(1982,8.54)(1983,4.61)(1984,5.01)(1985,6.01)
  (1986,3.42)(1987,4.18)(1988,4.93)(1989,7.80)(1990,9.47)
  (1991,5.87)(1992,3.57)(1993,1.60)(1994,2.42)(1995,3.48)
  (1996,2.40)(1997,3.18)(1998,3.40)(1999,1.51)(2000,2.98)
  (2001,1.75)(2002,1.67)(2003,2.90)(2004,2.99)(2005,2.83)
}; \addlegendentry{Gran Breta\~na}
\end{axis}
\end{tikzpicture}
\caption{Tasas de inflaci\'on anuales: los seis pa\'ises vs.\ EE.\,UU.\ (1981--2005).
  La l\'inea negra discontinua es la referencia de EE.\,UU.}
\end{figure}

% ---------------------------------------------------------------
\section{Comportamiento comparativo respecto a EE.\,UU.}
% ---------------------------------------------------------------

\begin{obs}
  \textbf{Co-movimiento general en los a\~nos ochenta.} Todos los pa\'ises
  registraron tasas de inflaci\'on elevadas a principios de los ochenta
  y las redujeron progresivamente a lo largo de la d\'ecada. La trayectoria
  descendente es com\'un, aunque los niveles difieren: Italia parte de
  $18.15\%$ en 1981, frente a $10.32\%$ de EE.\,UU.; Alemania, de $6.34\%$.
\end{obs}

\begin{obs}
  \textbf{Diferencial de nivel: Europa alta vs.\ Alemania baja.}
  Durante el per\'iodo completo, Italia y Gran Breta\~na mantuvieron
  tasas sistem\'aticamente superiores a las de EE.\,UU., mientras que
  Alemania operaba por debajo. Este patr\'on refleja diferencias en la
  credibilidad antiinflacionaria de los bancos centrales: el Bundesbank
  alemn goz\'o de la mayor reputaci\'on entre los pa\'ises del grupo.
\end{obs}

\begin{obs}
  \textbf{La anomal\'ia japonesa: desinflaci\'on hasta la deflaci\'on.}
  Jap\'on se desvincula del grupo a partir de 1995. Mientras EE.\,UU.\
  y Europa mantienen inflaciones positivas de entre $1\%$ y $4\%$, Jap\'on
  entra en territorio negativo (deflaci\'on) en 1999--2002 y en 2005.
  Este fen\'omeno corresponde a la trampa de liquidez y la d\'ebil demanda
  interna que caracterizaron a la econom\'ia japonesa tras el colapso de
  la burbuja de activos de 1989--1990.
\end{obs}

\begin{obs}
  \textbf{Convergencia en los noventa.} A partir de 1993--1994, las tasas
  de todos los pa\'ises (excepto Jap\'on) convergen hacia un rango
  $1\%$--$4\%$, pr\'acticamente indistinguible de las de EE.\,UU.\ Esta
  convergencia se acelera con la adopci\'on de metas de inflaci\'on
  (\textit{inflation targeting}) como r\'egimen de pol\'itica monetaria
  en el Reino Unido (1992), Canad\'a (1991) y Francia/Alemania (dentro
  del proceso de convergencia del euro, Maastricht 1992).
\end{obs}

% ---------------------------------------------------------------
\section{Causalidad vs.\ correlaci\'on: \textquestiondown{}la inflaci\'on
  de EE.\,UU.\ ``provoca'' la de los dem\'as?}
% ---------------------------------------------------------------

\subsection{El argumento correlacional y su l\'imite}

El hecho de que las tasas de inflaci\'on se muevan en la misma direcci\'on
no es evidencia de causalidad. En econometr\'ia, la confusi\'on entre
correlaci\'on y causalidad se conoce como la falacia
\textit{post hoc, ergo propter hoc}: el que un evento ocurra antes o junto
con otro no implica que lo cause.

\subsection{Causas alternativas del co-movimiento}

\begin{enumerate}[leftmargin=1.8em, itemsep=6pt]
  \item \textbf{Factores comunes externos.} Los dos grandes episodios de
    inflaci\'on alta (1973--1975 y 1979--1981) coinciden con los choques
    de oferta de la OPEP, que afectaron a todos los pa\'ises importadores
    de petr\'oleo de forma expl\'icitamente paralela y sin intermediaci\'on
    de EE.\,UU. La inflaci\'on no fue exportada por EE.\,UU.;\ fue
    importada por todos desde el mercado mundial del crudo.

  \item \textbf{Integraci\'on de mercados financieros.} Los tipos de inter\'es
    mundiales est\'an interconectados. Cuando la Reserva Federal elev\'o
    tasas a principios de los ochenta para combatir la inflaci\'on, los
    bancos centrales de Europa y Jap\'on respondieron de forma parecida
    para evitar depreciar sus monedas y, con ello, importar m\'as inflaci\'on
    v\'ia tipo de cambio. La simultaneidad no implica que EE.\,UU.\ dictara
    la pol\'itica; implica que los pa\'ises respondieron al mismo incentivo
    estructural.

  \item \textbf{Comercio internacional y transmisi\'on de precios.} Los precios
    de bienes transables se determinan en mercados globales. Una ca\'ida del
    precio del petr\'oleo o de las materias primas reduce la inflaci\'on en
    todos los pa\'ises simult\'aneamente, sin que medie un mecanismo causal
    que pase por EE.\,UU.
\end{enumerate}

\subsection{El test de Granger: condici\'on necesaria de causalidad}

Para afirmar que la inflaci\'on de EE.\,UU.\ \emph{causa} (en el sentido
de Granger) la de otro pa\'is $j$, deber\'ia cumplirse que la informaci\'on
pasada de $\pi_t^{\text{US}}$ mejore significativamente la predicci\'on
de $\pi_t^j$ por encima de lo que predice la historia propia de $\pi_t^j$:
\begin{equation}
  \pi_t^j = \sum_{k=1}^p \alpha_k \pi_{t-k}^j
          + \sum_{k=1}^p \beta_k \pi_{t-k}^{\text{US}}
          + \varepsilon_t
  \label{eq:granger}
\end{equation}
La hip\'otesis nula de ``no-causalidad de Granger'' es $H_0: \beta_1 = \cdots = \beta_p = 0$.
El simple co-movimiento que se observa en las gr\'aficas no proporciona
ninguna evidencia sobre esta hip\'otesis.

\begin{clave}
  \textbf{Respuesta a la pregunta del ejercicio.}\quad
  No. El co-movimiento de las tasas de inflaci\'on en la misma direcci\'on
  \emph{no indica} que la inflaci\'on en EE.\,UU.\ cause la de los dem\'as
  pa\'ises. Las razones son tres:

  \begin{enumerate}[leftmargin=1.4em, itemsep=4pt]
    \item Los movimientos paralelos pueden deberse a \textbf{factores
      comunes externos} (choques petroleros, ciclos de materias primas)
      que afectan a todos los pa\'ises simult\'aneamente.
    \item La \textbf{causalidad requiere precedencia temporal}: habr\'ia
      que demostrar que los cambios en la inflaci\'on de EE.\,UU.\
      \emph{preceden y predicen} los del pa\'is en cuesti\'on, una vez
      controlada la historia propia de ese pa\'is (test de Granger,
      ecuaci\'on~\ref{eq:granger}).
    \item La relaci\'on puede ser de \textbf{interdependencia rec\'iproca}:
      si EE.\,UU.\ eleva tasas de inter\'es para combatir su inflaci\'on,
      los dem\'as pa\'ises responden por razones propias (tipo de cambio,
      flujos de capital), no porque EE.\,UU.\ les ``transmita'' inflaci\'on
      directamente.
  \end{enumerate}
\end{clave}

\appendix

\section{Correlaci\'on con EE.\,UU.: estad\'isticos resumen}

\begin{center}
\begin{tabular}{lrrr}
\toprule
\textbf{Pa\'is} & $r(\pi^j, \pi^{\text{US}})$ & $r(\Delta\pi^j, \Delta\pi^{\text{US}})$
  & \textbf{Signo comun (\%)} \\
\midrule
Canad\'a    & 0.947 & 0.741 & 88 \\
Jap\'on     & 0.864 & 0.503 & 72 \\
Francia     & 0.981 & 0.795 & 84 \\
Alemania    & 0.860 & 0.612 & 76 \\
Italia      & 0.958 & 0.681 & 80 \\
Gran Breta\~na & 0.916 & 0.554 & 76 \\
\bottomrule
\end{tabular}
\end{center}

\noindent La columna $r(\pi^j, \pi^{\text{US}})$ es la correlaci\'on entre
niveles (muy alta, compatible con tendencia com\'un espuria); la columna
$r(\Delta\pi^j, \Delta\pi^{\text{US}})$ es la correlaci\'on entre las
primeras diferencias (m\'as baja, sugiere que parte de la correlaci\'on
en niveles es artefacto de la tendencia). ``Signo com\'un'' indica el
porcentaje de a\~nos en que ambas tasas cambiaron en la misma direcci\'on.

\vspace{14pt}
\noindent\rule{\linewidth}{0.4pt}

\noindent{\small\textit{Nota.} Los valores de correlaci\'on del Ap\'endice
son aproximaciones calculadas a partir de los datos de la Tabla~1.3.
El an\'alisis formal de causalidad de Granger requiere contrastar la
estacionariedad de las series antes de estimar \eqref{eq:granger}; v\'ease
el cap\'itulo~21 de Gujarati \& Porter.}

\end{document}
